\section{Opposite risk orderings after oracle tuning}
\label{sec:conference-predictive}

We compare attainable rare-presence risk after separate scalar retuning for
independent features ($\ell = 0$), $B = 1$, and the population preprocessing,
zero initialization, and fixed intercept of
Section~\ref{sec:conference-setting}. Put $b = q(1-q)$. Parameterize the
original-coordinate update matrices by
\begin{equation}
 D_U(\alpha) = \alpha I_2,\qquad
 D_S(\alpha) = \diag(4b\alpha,\alpha),\quad \alpha\geq0.
 \label{eq:conference-oracle-rates}
\end{equation}
Here $D=\eta P$, so $\eta_U=\alpha$ and $\eta_S=b\alpha$. Each
parameterization covers every nonnegative constant learning rate for its rule.
The shared parameter is a coordinate for the search, not a constraint that the
selected rates match.

For an observation cap $M\in\mathbb N$, define
\begin{equation}
 R_A^*(M)=\inf_{\substack{\alpha\geq0\\ k\in\{0,\ldots,M\}}}
             R_{R,k}^{A}(\alpha),\qquad A\in\{U,S\}.
 \label{eq:conference-oracle-risk}
\end{equation}
Rates and stops are deterministic oracle choices, based on the population risk
averaged over training randomness. Initialization is eligible. The search
includes finite stops at rates that need not be mean-square stable.

\begin{theorem}[Opposite oracle rare-risk orderings]
 \label{thm:conference-predictive}
 Suppose $\sigma^2>0$, $0<q\leq10^{-4}$, and
 $1\leq Mq(1-q)\leq4$ in the independent binary model above.
 \begin{enumerate}
  \item If $(\beta_R^*)^2=\sigma^2$ and
        $(\beta_C^*)^2\leq\sigma^2/100$, then
        $R_U^*(M)-R_S^*(M)\geq7\sigma^2/1000$.
  \item If $(\beta_C^*)^2=\sigma^2$ and
        $(\beta_R^*)^2\leq\sigma^2/100$, then
        $R_S^*(M)-R_U^*(M)\geq\sigma^2/200$.
 \end{enumerate}
 Either sign is allowed for each coefficient. Both inequalities also hold if
 each rule searches only its mean-square-stable nonnegative constant rates,
 together with initialization.
\end{theorem}

The proof in Appendix~\ref{sec:predictive-proofs} bounds every rate and stop of
the losing rule and supplies one stable candidate for the winning rule. The
diagonal error moments obey an exact positive recurrence, even though stochastic
updates couple the independent features. For a dominant rare signal, learning
it with U also incurs common-coordinate noise; S can reduce that cost. For a
dominant common signal, S must balance incomplete common learning against rare
noise. U can learn the common coordinate while leaving the small rare signal
largely unchanged. The bounds establish strict orderings over whole families of
procedures without identifying their optimizers.

The expected rare training count at the cap is $Mq=\tau/(1-q)$, where
$\tau=Mq(1-q)\in[1,4]$. Thus the positive gaps concern a regime with roughly
one to four rare observations, even as $M$ grows and $q$ decreases. Prevalence
alone does not determine the preferred normalization. The two signal families
give sufficient conditions, without locating a general transition between them.
The theorem does not optimize batches, learning-rate schedules, or
data-dependent stopping, and supplies no guarantee for estimated preprocessing
or validation selection.
